\section{结论}\label{sec:conclusion}

本文综述了大数定律从《猜度术》到当代的完整脉络。概括而言：

\begin{enumerate}[label=(\arabic*)]
    \item \textbf{历史主线}是“条件弱化、结论强化、刻画精确”：伯努利（1713）的频率稳定性经由切比雪夫的方差方法（1866）与马尔可夫的相依推广（1906），在博雷尔与康泰利（1909/1917）手中诞生强形式，由科尔莫哥洛夫（1928--1933）以极大值不等式、截断法与三级数定理完成独立情形的最终刻画，并经伯克霍夫（1931）推广到平稳遍历过程。
    \item \textbf{方法层面}，从切比雪夫不等式的方差控制，到博雷尔--康泰利引理的子序列法、科尔莫哥洛夫极大值不等式的路径控制、特征函数法、鞅方法与遍历论视角，每一次技术进步都对应着对“从单时刻估计到整条轨道控制”这一核心困难的新回答。
    \item \textbf{应用层面}，大数定律是统计相合性、蒙特卡洛方法、信息论典型性、遍历数论与精算科学的第一原理；它与中心极限定理、迭代对数定律构成描述部分和行为的完整“尺度阶梯”。
    \item \textbf{前沿方向}包括无穷均值情形的非经典规范化、强逼近与定量速率、相依框架的最优条件、巴拿赫空间与算子水平的无穷维推广，以及随机算法与量子概率中的结构化推广。
\end{enumerate}

作为概率论的第一极限定理，大数定律的历史恰是这门学科自身历史的缩影：从赌博问题的经验直觉，经分析不等式的锻造，成长为横跨纯粹数学与应用科学的普适原理。它所确立的信念——\textbf{随机性在大量重复之下显露出确定性的规律}——至今仍驱动着极限理论的每一次推广，也继续为数据时代的推断与计算提供着最根本的合法性证书。
